% !TEX root = ../main.tex
%% Sujet type bac n° 8 — sujet inédit OptiTAB (calculs vérifiés par programme : bac_s8.py)

\sujetBac{8}{QCM : probabilités et dénombrement/4,Suites et récurrence/5,Fonction exponentielle et convexité/6,Géométrie : distance à une droite/5}

\begin{exoBac}{4}{Questionnaire à choix multiples}
Pour chaque question, une seule des quatre réponses proposées est exacte. Aucune justification n'est demandée. Une réponse fausse, une
réponse multiple ou l'absence de réponse ne rapporte ni n'enlève aucun point.
\begin{questions}
  \item Le nombre d'anagrammes du mot LYCEE (mots de cinq lettres, ayant un sens ou non, formés avec ces cinq lettres) est :
\begin{tasks}[label=\alph*)](4)
  \task $120$
  \task $60$
  \task $25$
  \task $3\,125$
\end{tasks}
  \item $A$ et $B$ sont deux événements indépendants tels que $P(A)=0{,}4$ et $P(B)=0{,}5$. Alors $P(A\cup B)$ est égale à :
\begin{tasks}[label=\alph*)](4)
  \task $0{,}9$
  \task $0{,}2$
  \task $0{,}7$
  \task $0{,}45$
\end{tasks}
  \item Une variable aléatoire $X$ suit une loi binomiale d'espérance $12$ et de variance $3$. Ses paramètres sont :
\begin{tasks}[label=\alph*)](2)
  \task $n=24$ et $p=0{,}5$
  \task $n=48$ et $p=0{,}25$
  \task $n=12$ et $p=1$
  \task $n=16$ et $p=0{,}75$
\end{tasks}
  \item Une urne contient $10$ boules numérotées. Le nombre de façons de tirer simultanément $3$ boules de l'urne est :
\begin{tasks}[label=\alph*)](4)
  \task $720$
  \task $120$
  \task $1\,000$
  \task $30$
\end{tasks}
\end{questions}
\end{exoBac}

\begin{exoBac}{5}{Suites et récurrence}
On considère la suite $(u_n)$ définie par $u_0=1$ et, pour tout entier naturel $n$, $u_{n+1}=\dfrac12u_n+n+1$.
\begin{questions}
  \item Calculer $u_1$, $u_2$ et $u_3$.
  \item Démontrer par récurrence que, pour tout entier naturel $n$, $u_n\geq n$. En déduire la limite de la suite $(u_n)$.
  \item On pose, pour tout entier naturel $n$, $v_n=u_n-2n+2$. Démontrer que la suite $(v_n)$ est géométrique ; préciser sa raison et
    son premier terme.
  \item En déduire que, pour tout entier naturel $n$, $u_n=3\times\left(\frac12\right)^{n}+2n-2$.
  \item Déterminer la limite de la suite $\left(\dfrac{u_n}{n}\right)_{n\geq1}$.
  \item Écrire une fonction Python \code{somme(n)} qui renvoie $S_n=u_0+u_1+\dots+u_n$, puis donner une valeur approchée au millième de
    $S_{10}$.
\end{questions}
\end{exoBac}

\begin{exoBac}{6}{Fonction exponentielle et convexité}
On considère la fonction $f$ définie sur $\R$ par $f(x)=\left(x^{2}-3\right)e^{x}$, et l'on note $\mathcal{C}_f$ sa courbe représentative.
\partieExo{A}{Étude de la fonction}
\begin{questions}
  \item Déterminer les limites de $f$ en $-\infty$ et en $+\infty$. Interpréter graphiquement la limite en $-\infty$.
  \item Montrer que $f'(x)=(x-1)(x+3)e^{x}$ et dresser le tableau de variations de $f$, avec les valeurs exactes des extremums.
  \item Résoudre l'équation $f(x)=0$, puis étudier le signe de $f$.
\end{questions}
\partieExo{B}{Convexité et aire}
\begin{questions}[resume]
  \item On admet que $f''(x)=\left(x^{2}+4x-1\right)e^{x}$. Étudier la convexité de $f$ et donner les abscisses des points d'inflexion de
    $\mathcal{C}_f$.
  \item Vérifier que la fonction $F$ définie sur $\R$ par $F(x)=\left(x^{2}-2x-1\right)e^{x}$ est une primitive de $f$.
  \item Calculer l'aire, en unités d'aire, du domaine délimité par $\mathcal{C}_f$ et l'axe des abscisses, entre ses deux points
    d'intersection. On en donnera la valeur exacte puis une valeur approchée au centième.
\end{questions}
\end{exoBac}

\begin{exoBac}{5}{Géométrie : distance à une droite}
L'espace est muni d'un repère orthonormé. On considère la droite $d$ passant par $A(1\pv0\pv2)$ et de vecteur directeur
$\vec u\,(1\pv1\pv-1)$, ainsi que le point $M(3\pv4\pv1)$.
\begin{questions}
  \item Déterminer une représentation paramétrique de $d$. On note $M_t$ le point de $d$ de paramètre $t$.
  \item Montrer que, pour tout réel $t$, $MM_t^{2}=3t^{2}-14t+21$.
  \item Déterminer la valeur de $t$ pour laquelle la distance $MM_t$ est minimale. On note $H$ le point correspondant : donner ses
    coordonnées, ainsi que la distance minimale.
  \item Vérifier que la droite $(MH)$ est orthogonale à la droite $d$.
  \item Déterminer une équation du plan $\mathcal{Q}$ passant par $M$ et orthogonal à $d$, puis vérifier que $H$ est le point d'intersection
    de $\mathcal{Q}$ et de $d$.
  \item Déterminer les coordonnées des points de la droite $d$ situés à la distance $\sqrt6$ du point $M$.
\end{questions}
\end{exoBac}
